\documentclass[12pt]{article}

\usepackage[explicit]{titlesec}

\titleformat{\section}{\normalfont\Large\bfseries}{}{0em}{\thesection. #1}

% \usepackage[active,tightpage]{preview}
% \renewcommand{\PreviewBorder}{1in}

\title{The Northern Steppes Bylaws}
% \date{December 27, 2022}

\usepackage{enumerate}
\usepackage{enumitem}
\usepackage[margin=1in]{geometry}
\usepackage{xcolor}
\usepackage{color,soul}
% \usepackage[margin=0.75in, right=3.25in, marginparwidth=2.5in, marginparsep=0.25in]{geometry}
\usepackage[markup=bfit, addedmarkup=it, deletedmarkup=sout, commentmarkup=footnote, authormarkup=none, defaultcolor=red]{changes}
% \usepackage[markup=bfit, addedmarkup=it, deletedmarkup=sout, commentmarkup=margin, authormarkup=none, todonotes={textwidth=2.5in}]{changes}
% \usepackage{marginnote}

\usepackage[skip=1pt plus1pt, indent=40pt]{parskip}
% \setcommentmarkup{\marginpar{#1}}

\newlist{level}{enumerate}{10}
\setlist[level]{label*=\arabic*.}
\setlist[level,1]{label = \thesection.\arabic*.,ref = \emph{\thesection.\arabic*.}}

% \newcommand{\headItem}[1]{\global\isHeadtrue\item {\bfseries #1}}
% \newcommand{\maybeHeadItem}{%
%     \ifisHead\bfseries\global\isHeadfalse\fi%
% }
% \newif\ifisHead

\NewDocumentCommand{\subject}{ m }{%
    \begin{center}
        \bfseries%
        Re: #1%
    \end{center}
}%
\NewDocumentCommand{\name}{ m }{{\bfseries #1}}
\NewDocumentCommand{\email}{ m }{ #1 }

% \renewcommand{\thesection}{Article \arabic{section}}
% \renewcommand{\thesection}{Article \Roman{section}}

% \newcommand{\Newpage}{\end{preview}\begin{preview}}
% \newcommand{\newpart}[2][]{\colorbox{green!40}{#2}\expandafter\ifx\expandafter\relax\detokenize{#1}\relax\else\textnormal{\footnote{#1}}\fi}
\newcommand{\newpart}[2][]{\hl{#2}\expandafter\ifx\expandafter\relax\detokenize{#1}\relax\else\textnormal{\footnote{#1}}\fi}

\def\ex.#1\par{\section{#1}}
\def\a.#1{\begin{level}\item #1}
\def\b.#1{\item #1}
\def\z.{\end{level}}

\begin{document}
% \begin{preview}

% \maketitle

% \name{}

% \listofchanges

\ex. SAFETY AND ADJUDICATION:

    \begin{level}[start=0] \item 1.0. Each individual is expected to treat each other with respect and participate in fair play with safety and sportsmanship in mind. Actual violence, threats of actual violence, cheating, and arguing with marshals or other combatants disrespectfully are just a few examples of unacceptable behavior and are not allowed.
    \b. 1.1. Marshals: Marshals are the referees of Belegarth, responsible for rules enforcement, encouraging acceptable behavior and the inspection of both the field and equipment used.
        \a. 1.1.1. A marshal has the authority to remove anyone from the field of battle for reasons listed in rule 1.0.
        \b. 1.1.2. Marshals may call hit determinations at their discretion.
        \b. 1.1.3. Marshals oversee the safe conduct of battles and therefore have the power to declare “hold” whenever a safety concern arises.
        \b. 1.1.4. Marshals enforce the rules enumerated in Section 4 (Weapons Specifications and Checking), and determine the classifications of equipment.
        \b. 1.1.5. Marshals shall not be used in such a way as to gain protection or advantage in combat.
        \b. 1.1.6. Solid yellow tabards, baldrics, or other pieces of solid yellow garb signify that a person is a marshal or a non-combatant. Combatants are forbidden from wearing such.
        \b. 1.1.7. Non-combatants on the field must be approved by a marshal (e.g. photographers.)
        \b. 1.1.8. Intentional illegal acts may result in immediate removal from the field, until a marshal allows them re-entry. Marshals are the arbiter of intent.
        \b. 1.1.9. Unintentional illegal acts that re-occur frequently shall be treated as intentional.
        \z.
    \b. 1.2. Hold: In the event of an unsafe situation, anybody present MUST call “HOLD!” as loudly as possible. All activity must cease while a marshal assesses the situation. The battle will resume when the marshal calls “Positions” to prepare combatants for resuming and then “Lay on”.
        \a. 1.2.1. Combatants must drop to a knee immediately when a hold is called and call “HOLD!” as well. They should not move or re-equip during a hold unless told to do so by a marshal.
        \b. 1.2.2. It is illegal to call a hold or use a hold called by someone else to gain an advantage over fellow combatants.
        \b. 1.2.3. Combatants must avoid using phrases such as “Hold the line!” or any words that could be misinterpreted as “HOLD!” during a battle.
        \z.
    \b. 1.3. All equipment must be inspected and properly marked according to the guidelines outlined in Section 4.4 before it is used in combat.
    \b. 1.4. The target of an attack has the authority to make combat hit determinations, but must defer to the discretion of marshals.
    \b. 1.5. Creative interpretation of the rules to gain any advantage is illegal. The marshal settles all disputes.
    \b. 1.6. Fighting near weapon piles, onlookers or unsafe locations is discouraged.
    \b. 1.7. A combatant may choose to call themselves dead and/or leave the field at any time by displaying death in an appropriate manner, rule 3.7.3.4.
    \z.



1.7.1. Combatants may not re-enter the field until game-play has ceased or a marshal has allowed re-entry. This must be done in a safe manner to get clear of the battlefield boundaries and should not interfere with the ongoing battle.

2. EQUIPMENT CLASSIFICATIONS AND DEFINITIONS:
2.1. There are five classifications of offensive equipment, hereafter called weapons. Weapons are any items that can score combat hits. All weapons must meet the requirements outlined in Section 4, Weapon Specifications and Checking.

2.1.1. Class 1: One-handed swung weapon.
2.1.2. Class 2: Two-handed swung weapon.
2.1.3. Class 3: Thrusting weapon.
2.1.4. Class 4: Missile weapon.
2.1.5. Class 5: Head-only missile weapon.

2.2. Defensive equipment is any item that gives combat advantage to its wielder by preventing combat hits and is unable to inflict damage on combatants.

2.2.1. There are two types of defensive equipment: shields and armor. All defensive equipment must meet the requirements outlined in section 4.11 and 4.12, respectively.
2.2.2. Shields are rigid objects that are padded on the front and sides and are equipped with handles or straps.
2.2.3. A shield may not be constructed in a manner that would confer the advantage of unbreakable armor.

2.2.3.1. Unbreakable armor is defined as a worn shield that shapes/forms closely to more than one target area, wraps around a limb unreasonably, and/or shapes/forms closely to multiple planes of the torso.

2.2.4. Armor is protective body covering, consisting of period materials. Armor must be readily recognizable as armor.

2.3. No single piece of equipment can be classified as both offensive and defensive equipment, e.g., a weapon cannot also serve as a shield or armor.
2.4. Miscellaneous equipment includes, but is not limited to, items such as: thin belts, pouches, boots, quivers, straps, scabbards, and non-armor clothing and headwear.

2.4.1. While conferring no special rules advantage, miscellaneous equipment may be checked for combat safety and Belegarth-appropriate appearance at the marshal’s discretion. The minimum non-armor clothing requirements are outlined in Section 5, (Garb).
2.4.2. Miscellaneous equipment cannot serve as offensive or defensive equipment regardless of the material(s) used.

3. COMBAT:
3.1. The following actions are allowable contact:

3.1.1. Weapon to weapon contact.
3.1.2. Weapon to body contact is allowed on valid target areas using striking-legal surfaces.
3.1.3. Body to weapon contact: Pushing, grabbing, or sweeping a combatant’s strike-legal surface results in a valid hit to the body location used for contact. Additional rules may be found in 3.12.
3.1.4. Weapon to shield contact is allowed.
3.1.5. Shield to weapon contact: Shields may be used to strike, deflect, move or pin a combatant’s weapon.
3.1.6. Shield to shield contact: Shields may be used to strike, deflect, move, or pin a combatant’s shield.
3.1.7. Shield to body contact: Shields may be used to strike, deflect, move, or pin a combatant’s body other than what is mentioned in 3.2.
3.1.8. Body to shield contact: Combatants may manipulate another combatant’s shield with their body, including feet, knees, shoulders, hands, and elbows.
3.1.9. Body to Body contact: Combatants may make contact with other combatants in accordance to rule 3.12.

3.2. Disallowed contact:

3.2.1. Head contact from Class 1, 2, 3 weapons or shields.

3.2.1.1. Feints towards the head from Class 1, 2, 3 weapons or shields are discouraged.

3.2.2. Unarmed punches and kicks directed at other combatants. bodies.
3.2.3. Throws where the throwing combatant allows the other combatant to freefall to the ground.
3.2.4. Intentionally hitting another combatant with a non-striking surface.
3.2.5. Joint/nerve holds and manipulations.
3.2.6. Grappling the head or neck.

3.3. Target area definitions: Hits

3.3.1. Body: Area bounded by the base of neck (inclusive), shoulder-arm joint (inclusive), hip-leg socket (inclusive), groin, and buttocks (inclusive).
3.3.2. Arm(s): Area bounded by the wrist (inclusive) and the shoulder-arm joint (exclusive).
3.3.3. Leg(s): Area bounded by the ankle (inclusive) and hip-leg socket (exclusive).
3.3.4. Head: Area above the base of neck (exclusive).
3.3.5. Hand(s): Area below the wrist (exclusive). An empty hand is a legal target area. A hand on a weapon or shield is considered part of that weapon or shield. Any hit to the hand is considered a hit to the arm.
3.3.6. Feet: Area below the ankle (exclusive). A foot is a legal target area if it is off the ground. Any hit to the foot is considered a hit to the leg.

See visual guide on Belegarth hit locations.
3.4. Weapons:

3.4.1. Weapons which strike with sufficient force can score a hit to the target area.
3.4.2. “Sufficient force” is defined as being both “solid” and having powerful impact on a target area as defined below:

3.4.2.1. Solid: Successfully strikes the target area. Taps, grazes, significantly obstructed strikes, and strikes that come in contact with garb only do not count as sufficient force.
3.4.2.2. Strikes that do not have sufficient force shall be communicated as insufficient by saying “light”, “graze”, or “garb’ as appropriate.

3.5. Weapon Damage: Weapons yield various amounts of damage according to the classification of the weapon and the armor/damage status of the target.

3.5.1. Class 1 (one-handed, swung) weapons cause one hit to a target area. Any Class 1 or Class 2 weapon swung with one hand, no matter the length, is a Class 1 weapon. Class 1 weapons swung with two hands causes one hit to a target area.
3.5.2. Class 2 (two-handed, swung) weapons cause two hits to the target area when used for two-handed strikes.

3.5.2.1. “Two-handed” is defined as having both hands fixed on the weapon when the weapon makes contact with the target.
3.5.2.2. Combatants striking with a two handed swing should call “Two” or the appropriate tape color as indicated on the weapon as they strike.

3.5.3. Class 3 (thrusting) weapons wielded one-handed cause one hit of damage to an unarmored target area, and have no affect against an armored area.

3.5.3.1. When used two handed, Class 3 weapons bypass armor.
3.5.3.2. Combatants striking with a two handed stab should call “Double” as they strike. See 3.5.2.1. For two-handed definition.

3.5.4. Class 4 (missile) weapons cause one hit to a target area and bypass all armor except head armor. A Class 4 weapon striking an armored portion of the head area causes no hit.
3.5.5. Class 5 (head only missile) weapons cause 1 hit to an unarmored head area. A Class 5 weapon striking an armored portion of the head area causes no hit.
3.5.6. The head is a legal target area for Arrows/Bolts, Thrown Javelins, and all Class 5 weapons

3.5.6.1. The head is an illegal target area for Class 1, 2, and 3 weapons.

3.5.7. A Class 1 or Class 2 weapon cannot also be a Class 4 or 5 weapon.
3.5.8. Intentionally hitting combatants with the non-striking surface of a weapon is illegal. For example, deliberately striking a combatant with flail haft padding to allow the head to swing around and hit a target area.
3.5.9. Anvilling is blocking a weapon strike by laying a weapon against a target area and/or shield and is illegal.

3.5.9.1. Sufficient force hits must be taken through anvilling weapons.
3.5.9.2. Hits must be taken from weapons not in direct contact with a target area or shield but are driven into a target area or shield with sufficient force.

3.5.10. Sheathed or otherwise worn weapons cannot block attacks.
3.5.11. Shot in Motion: An attack on another combatant that begins before the attacker is hit by a legal blow is considered a legal shot on the combatant. The shot may not change planes once the attacker is hit (e.g. a fake on the combatant).
3.5.12. Magic Switch: It is permissible to move a weapon or shield from the disabled hand to the uninjured hand immediately after a combatant’s arm is hit.

3.6. Armor:

3.6.1. Armor confers one additional hit to the target area covered by the armor. Multiple pieces of armor on the same target area only confer a single hit. A single piece of armor covering multiple areas confer a hit on each target area covered.
3.6.2. Armor only protects areas covered.
3.6.3. Armor must cover at least one-third of a target area. Armor which extends continuously from a different target area is not required to cover one-third of a target area to count as armor for that target area.
3.6.4. Weapons that strike both armored and unarmored target areas are considered to have hit the unarmored target area.
3.6.5. The presence of armor must be easily discernible to count as armor.
3.6.6. “Armor” must be declared to acknowledge that a sufficient-force strike hit the armor but did not disable a target area.

3.6.6.1. It is encouraged to include the target area in the declaration of armor; e.g. “left leg armor”, or “body armor”.

3.7. Hits:

3.7.1. All hits and armor status effects must be accurately portrayed at all times and truthfully reported when asked.
3.7.2. Effects of hits:

3.7.2.1. One hit to an unarmored target area disables that target area.
3.7.2.2. Two hits to an armored target area disable that target area.
3.7.2.3. Disabled arm: A disabled arm may not hold anything. If the arm is hit by a Class 1 or 2 weapon, the arm must be placed behind the back. If the arm is hit by a Class 3 or 4 weapon, leave the arm dangling limply.
3.7.2.4. Disabled leg: kneel on ground with the non-disabled leg up. When hit, the combatant must immediately drop to the disabled knee.

3.7.2.4.1. A combatant who has their leg disabled may either shuffle on the knees or move in any way that doesn’t involve weight on the foot of the disabled leg.
3.7.2.4.2. Knee shuffling is defined as moving on the knees such that one knee is on the ground at all times.

3.7.2.4.2.1. Examples of legal non-shuffling movements include crawling, rolling, lunging from the non-injured leg, or being clearly supported by others.
3.7.2.4.2.2. Examples of illegal movement include duck walking, lunging off the injured leg, or standing using the uninjured leg in any way.

3.7.2.4.3. If a combatant has both knees on the ground or both knees in the air, a strike to either leg is considered to have hit the good leg.
3.7.2.4.4. It is illegal to change your disabled leg, unless a medical condition requires you to do so.
3.7.2.4.5. Combatants must verbally distinguish between a leg disabled by a Class 1 or 2 weapon or a leg disabled by a Class 3 or 4 weapon.
3.7.2.4.6. A combatant who has had their leg disabled may choose to “post” instead of kneeling. When posting, the injured leg may not be moved, except to pivot in place. When posting, a strike to either leg is considered to have hit the good leg. When choosing to post, it must be declared at the time of posting and communicated normally as any other injury.

3.7.2.5. A disabled body causes death.
3.7.2.6. A disabled head causes death.
3.7.2.7. Two disabled limb target areas (arms and/or legs) cause death.

3.7.2.7.1. Limbs injured with Class 3 or Class 4 weapons do not count towards this total.

3.7.3. Death: Combatants must lay down immediately. Combatants are only allowed to move if instructed by a marshal or in order to move away from a potentially unsafe situation.

3.7.3.1. Attempting to gain a combat advantage over “living” players by appearing dead or declaring death and then returning to play is illegal. (e.g. sitting down appearing to look dead and waiting for someone to draw near.)
3.7.3.2. Combatants cannot return to life or otherwise undo a bad call unless otherwise specified by a marshal to do so.
3.7.3.3. Combatants may communicate “late” when their hit lands after their death, nullifying the damage from the strike to other combatants.
3.7.3.4. A combatant may indicate that they are dead by placing a weapon or arm on their head and loudly calling “DEAD”. This is only allowed when dead combatants are attempting to exit the field as instructed to by a marshal, or to call themselves dead, as in rule 1.7.
3.7.3.5. Dropping weapon(s) is not a valid show of death.
3.7.3.6. If a combatant is dead, they must look dead and make it clear to those around them that they are dead. (e.g. placing an elbow on the ground while lying down).
3.7.3.7. Dead combatants must not talk to the living unless to indicate a potential safety hazard.

3.7.4. Combatants attacking an unaware combatant with a Class 2 or 3 weapon must shout “Two” or appropriate tape color as indicated with the weapon they strike with a two-handed Class 2 swing, “Single” with a one-handed Class 3 shot, and “Double” if with two-handed Class 3 attack as appropriate. If the weapon class is not called, the combatant should consider a successful strike to cause a single hit.
3.7.5. Subsequent hits to the same location:

3.7.5.1. All subsequent hits with Class 3 or 4 weapon on the same target area previously injured only by a Class 3 or 4 weapon are ignored.
3.7.5.2. All subsequent hits to an arm disabled by a Class 1 or 2 weapon pass through to the body. However, armor still provides its protective benefits in the case of subsequent hits. For example, if a combatant has an arm disabled but is wearing torso armor, a subsequent Class 1 hit to the arm would first count against the armor and the following hit would be to the body.
3.7.5.3. All subsequent hits to a leg disabled by a Class 1 or 2 weapon during the initial movement of the knee to the ground constitute death. Once the knee is on the ground, subsequent shots to the disabled leg are ignored.
3.7.5.4. A target area disabled by a Class 3 or 4 weapon that is subsequently hit by a class 1 or 2 weapon is then considered to be disabled by a Class 1 or 2 weapon.

3.8. A hit that strikes both the body AND either an arm or a leg is assumed to have hit the body.
3.9. A single strike can only damage one target area.
3.10. Shields:

3.10.1. Shields are destroyed by two heavy, solid, two-handed strikes from a Class 2 weapon.

3.10.1.1. Subsequent strikes to a destroyed shield continue into the target area on which the shield is worn. For example, if a shield on an arm is broken by two sufficient Class 2 hits and is not dropped, the next hit would be to the arm and, after that, to the body.

3.10.2. Heavy strikes are defined as a stronger than normal strike, as defined in 3.4.2 and 3.5.2.1.
3.10.3. Shields may be used in any reasonable manner and still be considered a shield.
3.10.4. Only one shield may be used by a combatant at a time.
3.10.5. The wielder of the shield determines if a shield breaking hit is sufficient.
3.10.6. Shields lying on the ground cannot be broken.

3.11. Offensive Shield Techniques:

3.11.1. It is illegal to use offensive shield techniques to move combatants into hazards or obstacles such as trees or paved surfaces.
3.11.2. It is illegal to use the unpadded portions of the shield for offensive shield techniques.
3.11.3. Intentional shield contact to the head or neck is illegal.
3.11.4. Shield Bashing and Checking:

3.11.4.1. Shield Bashing is defined as a combatant charging another combatant and using the face or edge of their shield to make contact with their shield or body, such that forward momentum is impossible to stop within two steps.
3.11.4.2.Shield checking is defined as a combatant using the face or edge of their shield to make contact with another combatant while stationary or charging from two steps away or less, such that the combatant is able to stop their forward momentum within two steps.
3.11.4.3. It is legal to shield bash or check another combatant from the front or the sides excluding the rear quadrant. The combatant initiating the bash or check must ensure that this will not place another combatant in an unsafe situation.
3.11.4.4. A combatant may bash another combatant who does not have a shield.
3.11.4.5. While performing a shield bash or check, a combatant may not make intentional contact with another combatant’s head or neck.
3.11.4.6. Bashes must target a combatant’s center of mass, not knees or legs.
3.11.4.7. It is illegal to bash or check combatants that have bows and/or arrows/bolts.
3.11.4.8 It is illegal to shield bash, as defined by 3.11.4.1, a combatant with a disabled leg
3.11.4.9 It is legal to shield check, as defined by 3.11.4.2, a combatant with a disabled leg.

3.11.5. Shield braces, edging, and bumping:

3.11.5.1. A shield brace is when a combatant plants their feet while holding or placing their shield in front of a moving combatant.
3.11.5.2. A shield bump is incidental shield contact against a combatant’s body or equipment when the intent is not to knock the combatant over.
3.11.5.3. Edging is using the edge of a shield against a combatant’s shield, body, or weapons.
3.11.5.4. Shield braces and bumps are legal from all four sides against other combatants.
3.11.5.5. It is illegal to intentionally edge a combatant’s head or neck.
3.11.5.6. Edging is legal from the front and both sides but not from the rear.

3.11.6. Shield Kicking:

3.11.6.1. Shield Kicking is when a combatant makes contact with another combatant’s shield with their foot.
3.11.6.2. Kicking is allowed only to shields, not to people. The kicker must maintain sufficient control to not kick another combatant.
3.11.6.3. Kicking of shields less than 18. in diameter is illegal.
3.11.6.4. A shield kicker must maintain one foot on the ground. Kicks where both feet leave the ground are illegal.
3.11.6.5. Shield kicking from the rear is illegal.

3.12. Grappling:

3.12.1. Grappling is allowed.
3.12.2. Grappling is defined as wrestling or attempting to grasp another combatant’s body to prevent attacks or to control the combatant’s movements.
3.12.3. Grasping the non-striking surfaces of another combatant’s weapon such as a haft of a flail or a spear shaft does not constitute a grapple.
3.12.4. Combatants may initiate grapples with other combatants according to the following rules:

3.12.4.1. A combatant wearing no armor may grapple all combatants.
3.12.4.2. A combatant wearing leather or mostly (2/3 or more) leather composite armor may grapple any armored combatant, but not unarmored combatants.
3.12.4.3. A combatant wearing chain armor or mostly (2/3 or more) metal composite armor may grapple combatants wearing mostly metal composite, chain or plate armor.
3.12.4.4. A combatant wearing plate armor may not initiate a grapple.
3.12.4.5. A combatant wearing plastic safety equipment is treated as leather armor for grappling purposes only.
3.12.4.6. Groin protection, protective sports bra inserts, and safety glasses are exempt from 3.12.4.5.

3.12.5. A grappler must maintain positive control of the combatant when attempting to bring a grappled combatant to the ground.

3.12.5.1. Positive Control is defined as a grappler bearing some of a combatant’s weight and speed when bringing them to the ground.

3.12.6. Combatants with bows and/or arrows/bolts may not initiate grapples or be grappled.

3.12.6.1. Combatants may not grasp a combatant’s bow or arrows/bolts with the intent to control their direction or prevent them from using this equipment.

3.12.7. Combatants may grab their weapons any way they wish, including the blade/striking surface. This is an exemption to the anvilling rule (3.5.9).
         3.12.7.1. If holding the weapon from a striking surface, a combatant cannot attack to deal damage to other combatants in an unsafe manner (e.g. punching and/or gripping so that the handle or pommel is swung towards an opponent).
3.12.8. Gripping the striking surface of a combatant’s weapon results in the disabling of that limb.
3.12.9. Contested weapons, e.g., weapons being held by two or more combatants vying for control, may not be held by the striking surfaces. A combatant who grabs the striking surface of a contested weapon must lose that limb.

3.13. Missile Weapon Conventions:

3.13.1. If a bow or crossbow is hit by a Class 1 or 2 weapon, it is considered broken and cannot be used for the remainder of the battle.
3.13.2. A half draw or throw for Class 4 weapons under a range of 20 feet is required.

3.13.2.1. Half draw is defined as drawing back the bow only so far such that it imparts no more than half the force to the arrow than a normal full draw. This distance and pull will vary due to variances in bow design.
3.13.2.2. Half throw is defined as drawing back the javelin only so far such that it imparts no more than half the force of a normal full throw.

3.13.3. A missile weapon must travel its entire length to score a hit.
3.13.4. An arrow must be fired from a bow in order to score a combat hit. Similarly, a bolt must be fired from a crossbow in order to score a combat hit.
3.13.5. A missile weapon is considered to have hit if there is significant deflection of the missile head (greater than 30 degrees). Once the missile head has significantly deflected, the missile is rendered harmless until retrieved and fired again.
3.13.6. An archer who attacks with an arrow or bolt may call a combat hit for clarification when the shot clearly and unambiguously hit a target area.

3.13.6.1. For a shot to be clear and unambiguous, the archer must have an unobstructed view of the entire flight of the arrow or bolt including post hit deflection.

3.13.7. A javelin thrower may clarify that the javelin did not hit with the point by declaring “shaft” or saying that the shot was not good.
3.13.8. Javelin throwers may call “Point” for clarification.
3.13.9. When in doubt, the target makes the hit determination for missile weapons.
3.13.10. Blocking Missiles.

3.13.10.1. All rocks and javelins may be blocked by any means that keeps the missile away from a target area.
3.13.10.2. An arrow or bolt may only be blocked by a shield. An arrow or bolt blocked by a weapon is considered to have continued to travel in the same direction and strike the target area immediately behind the weapon.
3.13.10.3. Intentional blocking of an arrow or bolt with anything but a shield causes death to the blocker. This includes attempting to swat or grab arrows or bolts out of the air using weapons or limbs.

4. WEAPON SPECIFICATIONS AND CHECKING.
4.1. Definitions:

4.1.1. Striking surface: Padded surface of a weapon designed to make contact with a combatant during combat. Only the striking surface of a weapon may score a hit.
4.1.2. Non-striking surface: Any padded surface of the weapon that is not a striking surface.

4.1.2.1. Incidental Padding: The padded part of a weapon that could make direct contact with a combatant during a swing but cannot score a hit.
4.1.2.2. Courtesy Padding: The padded part of a weapon that might make contact with a combatant during combat, but is unlikely to do so in the direction of the swing. For example, the haft padding for a spear or just above a flail’s handle.

4.1.3. Handle: Non-padded portion of the weapon designed as a handhold.
4.1.4. Pommel: Non-striking surface that covers the end of the handle.
4.1.5. Crossguard: Non-striking surface that separates the striking surface from the handle and is perpendicular to the striking section of the weapon.
4.1.6. Hilt: The combination of the handle, pommel, and crossguard.
4.1.7. Core: The center of the weapon used to provide rigidity and flexibility to the striking and non-striking surfaces attached to it.
4.1.8. Sword: Any weapon approximating a medieval sword, constructed using either an edge/flat or cylindrical design.
4.1.9. Flail: Any hinged weapon.
4.1.10. Double-ended weapon: A weapon approximating a medieval staff.
4.1.11. Javelin: When thrown a Class 4 weapon, when used to stab a Class 3 weapon.
4.1.12. Archery: Class 4 weapons including bows, crossbows, arrows, and bolts.
4.1.13. Rocks: Class 5 weapons.

4.2. Padding and cover requirements:

4.2.1. Padding on the striking surface must have sufficient cushioning to prevent the core from being felt during a full-force hit.

4.2.1.1. Padding on the striking surface must have sufficient cushioning to prevent excessive stinging or bruising during a full-force swing.
4.2.1.2. All striking surfaces must be covered by cloth. The cloth covering must be in good condition with no serious rips or holes, especially near the tip of a weapon.
4.2.1.3. Cloth grip tape is not an allowable substitute for covering striking surfaces.

4.2.2. Incidental Padding is required on those parts of a weapon that are likely to come into contact with a combatant during normal combat but are not striking surfaces.

4.2.2.1. Incidental Padding must be sufficiently cushioning as to prevent the core from being felt during a sufficient force shot.
4.2.2.2. Incidental Padding is not required to be as soft as striking surface padding.
4.2.2.3. Incidental Padding must be covered by cloth, tape or plastidip.

4.2.3. Courtesy Padding may be used in any situation where padding is required, but it is unlikely that the weapon will contact a combatant.

4.2.3.1. Courtesy Padding must have sufficient cushioning to provide some protection in the event of contact.
4.2.3.2. Courtesy Padding must be covered by cloth, tape or plastidip.

4.2.4.Hafted: A weapon is hafted if the striking surface does not extend all the way to the handle.

4.2.4.1. Hafted Class 1 weapons must have incidental padding that extends at least 6 inches (15.24 cm) from the striking surface, or to the handle, whichever is shorter, and courtesy padding that extends from the bottom of the incidental padding to the handle.

4.2.4.2. Two-handed hafted Class 2 only and Class 2/Class 3 combined weapons must have incidental padding that extends at least 12 inches (30.48 cm) from the striking surface or to the handle, whichever is shorter, and courtesy padding that extends from the bottom of the incidental padding to the handle.

4.2.5. Handles must be contiguous. It is not allowable to have a handle, courtesy padding, and then an additional unpadded handle closer to the striking surface(s).
4.2.6. The only exception to this rule are double-ended weapons with specifications that conform to 4.9.6.

4.3. Tape may be used on striking surfaces under the cloth covering. However, the tape may not cause the weapon to hit too hard as determined by a hit test.
4.4. Marking: Weapons must be marked with the appropriate color(s) of tape to denote their classifications. This marking tape must be placed in a manner so that combatants and marshals may easily see it. Unmarked weapons will only be checked as Class 1 weapons.

4.4.1. Class 1 weapons are marked with blue tape on either the pommel or handle.
4.4.2. Class 2 weapons are marked with red tape on either the pommel or handle.
4.4.3. Class 3 weapons are marked with green tape on either the pommel or handle.

4.4.3.1. Class 1 and 2 weapons that are also Class 3 are marked in the same way.

4.4.4. Class 4 and 5 weapons are marked in a manner to indicate a marshal has inspected them.

4.5. Template rules:

4.5.1. Two and one-half inch rule (6.35 cm): No surface on a striking edge (sword tip, arrow/bolt head, spear head, javelin head, etc.) whether designed for stabbing or not, may readily pass more than one-half (1/2) inch (1.3 cm) through a two and one-half (2 1/2) inch (6.35 cm) hole.

4.5.1.1. The weapon tip is exempt from the two and one-half (2 1/2) inch (6.35 cm) rule, rule 4.5.1, if the weapon has a semicircular tip with a minimum one and one-half (1 ½) inch (3.81 cm) radius.

4.5.1.1.1. When checking a Class 1, Class 2, Class 3, or Class 4 weapon’s tip, the template should be applied perpendicular to the end of the weapon. (ie: place the template flat on top of the end of the weapon.)

4.5.2. Weapon pommels and other protrusions, such as the ends of crossguards, may not readily pass more than one-half (½) inch (1.3 cm) through a two (2) inch (5.08 cm) diameter hole.

4.6. The maximum allowed flex of any weapon except javelins is forty-five (45) degrees.

4.6.1. Arrows, bolts, bows, crossbows, and Class 5 weapons are exempt from flex rules.

4.7. All areas of wood-cored weapons must be taped, including bamboo and rattan.

4.7.1 Weapon handles and wooden bows are not required to be taped.
4.7.2. Wooden arrow/bolt shafts must be wrapped completely in tape prior to building the arrowhead.

4.8. Other than aluminum arrow/bolt shafts, a weapon may not have a metal core.
4.9. All weapons must be built to the following specifications:

4.9.1. Class 1: All Class 1 weapons must conform to the following, as applicable:

4.9.1.1. A Class 1 weapon under twenty-four (24) inches (60.96 cm) in length has no weight minimum.
4.9.1.2. A Class 1 weapon twenty-four (24) inches (60.96 cm) in length or longer must weigh a minimum of twelve (12) ounces (340.2 g).
4.9.1.3. With the exception of double-ended weapons, a Class 1 weapon must be shorter than forty-eight (48) inches (121.92 cm).
4.9.1.4. The maximum handle length for a Class 1 weapon is twelve (12) inches (30.5 cm) or one-third (1/3) of the overall length, whichever is greater. This cannot exceed one-half (1/2) of the overall length.
4.9.1.5. The minimum overall length of a Class 1 is twelve (12) inches (30.5 cm) of contiguous striking, incidental, and courtesy padding plus the length of the hilt.
4.9.1.6. A Class 1 weapon may also be Class 3.

4.9.2. Class 2: All Class 2 weapons must conform to the following:

4.9.2.1. The minimum length is forty-eight (48) inches (121.92 cm).
4.9.2.2. The minimum weight is twenty-four (24) ounces (680.4 g).
4.9.2.3. The maximum handle length for Class 2 weapons is eighteen (18) inches (45.72 cm) or one-third (1/3) of the overall length, whichever is greater. This cannot exceed one-half (½) of the overall length.
4.9.2.4. A Class 2 weapon may also be Class 3.

4.9.3. Class 3: All Class 3 weapons must conform to the following:

4.9.3.1. If the weapon is Class 3 only, it has no weight restriction.
4.9.3.2. The maximum handle length for Class 3 weapons is two-thirds (2/3) of its overall length.
4.9.3.3. If the weapon is Class 3 only, it may not have a yellow cover.

4.9.4. Single-edged weapons must have their non-striking edge clearly marked for at least twelve (12) inches (30.5 cm) with tape, paint, fabric, or other material in a way that contrasts with the striking surface cover and does not wrap onto the flat of the blade.


4.9.5. Flails must conform to the following:

4.9.5.1. The minimum striking surface circumference measured on the smallest plane passing through the center of the flail head is fifteen inches.
4.9.5.2. Only the head of a flail is a striking surface.
4.9.5.3. The maximum chain/hinge connecting the striking surface to the hilt is six (6) inches (15.24 cm).
4.9.5.4. Only one hinge per flail is allowed.
4.9.5.5. The hinged part of the flail must be padded with foam (often referred to as dingleberries) to keep the chain from easily entangling a weapon or body part. No more than one and one-half (1 ½) inches  (3.81 cm) of chain may be exposed.
4.9.5.6. The maximum overall length is forty (40) inches (101.6 cm).
4.9.5.7. Flails must contain incidental padding for the last six (6) inches (15.24 cm) of the haft before the chain.

4.9.6. Double-ended weapons must conform to all of the following:

4.9.6.1. Double-ended weapons must not be more than seven (7) feet long (2.13 m).
4.9.6.2. Double-ended weapons must have a minimum of eighteen (18) inches (45.72 cm) in length of striking surface covering each end in a cylindrical fashion. Both striking surfaces of this weapon must follow Class 3 weapon standards for a double-ended weapon to be legal.
4.9.6.3. Regardless of length, a double-ended weapon is a Class 1 weapon when swung and Class 3 when thrust.
4.9.6.4. Double-ended weapons may have no more than one-third (1/3) its overall length as unpadded handle.

4.9.7. Javelins must conform to all of the following:

4.9.7.1. Must also pass as a Class 3 weapon.
4.9.7.2. The maximum weight is sixteen (16) ounces (453.6 g).
4.9.7.3. The minimum length is four (4) feet (1.22 m).
4.9.7.4. The maximum length is seven (7) feet (2.13 m).
4.9.7.5. Must have courtesy padded along the entire length.
4.9.7.6. Must flex less than ninety (90) degrees. This is an exception to rule 4.6.
4.9.7.7. Must have a yellow cover on the striking surface of the weapon.

4.9.8. Archery Restrictions:

4.9.8.1. Arrows and bolts must conform to the stated arrow construction requirements and are exempt from non-arrow weapon construction requirements.
4.9.8.2. Compound bows or compound crossbows are not allowed.
4.9.8.3. Bows may not have any dangerous protrusions, such as metal post arrow rests, or mechanical modifications such as sights, stabilizers, or releases.
4.9.8.4. The maximum poundage allowed on a bow is thirty-five (35) lbs (15.88 kg) pull at twenty-eight (28) inches (71.12 cm) of draw.
4.9.8.5. A crossbow’s draw may not exceed 450 inch-pounds (518.5 kgf-cm).

A crossbow’s draw may not exceed 450 inch-pounds (518.5 kgf-cm). Inch pounds are calculated by multiplying ‘Draw Weight’ by ‘Power Stroke’. Draw Weight is measured at the bow’s full draw. Power Stroke is the distance from rest to full draw measured in inches. See Appendix C for reference chart.

4.9.8.6. A draw stop is required and must effectively stop an arrow from being drawn more than twenty-eight (28) inches (71.12 cm). It should protrude at least one-fourth (¼) of an inch (6.4 mm) away from the arrow shaft.

4.9.8.6.1. If the base of the head of an arrow prevents the archer from drawing beyond 28 inches (71.12 cm) the head of the arrow acts as the draw stop.

4.9.8.6.2. Crossbow bolts must have a maximum shaft length of no more than 18 inches (45.7 cm). A draw stop may not be used to meet this requirement. Crossbow bolts must use non-pinching nocks designed for use specifically with crossbows. Examples of these include half-moon, flat, capture, or multi grooved nocks.

4.9.8.7. Arrow/bolt striking surfaces may not easily pass more than one-half (½) inch (1.27 cm) through a two and one-half (2 ½) inch (6.35 cm) diameter hole. No part of the arrow/bolt’s striking surface may be less than two and one-half (2 ½) inches (6.35 cm) in any direction.
4.9.8.8. All arrows/bolts must contain a penny, or solid metal blunt of an equivalent gauge and circumference, perpendicularly secured at the end of the shaft.

4.9.8.8.1. All arrows/bolts using modular technology must create a semi-permanent connection point through the means of threaded screws, epoxy, glue, or strapping tape; the head must be secondarily secured at the end of the shaft with tape.
4.9.8.8.2. All arrows/bolts that are altered in any way during a day of combat will be treated as new arrows/bolts and must be rechecked as such before being put back into use.

4.9.8.9. The arrow’s/bolts striking surface must be constructed of open-cell foam.
4.9.8.10. All arrows/bolts must have at least two full fletchings.

         4.9.8.10.1. Fletchings must be made of a flexible material and cannot be rigid
4.9.8.11. The striking surface of an arrow/bolts must be free of tape.
4.9.8.12. The arrowhead should not have excess axial or lateral movement and must be secured at the end of the shaft in such a way that they will not come off if firmly twisted or firmly pulled.

4.9.9. Class 5 weapons have a minimum diameter of four (4) inches (10.16 cm) and are constructed entirely of foam, cloth and/or tape (coreless).

4.10. Prohibited Weapons:

4.10.1. Entangling weapons (nets, lassos).
4.10.2. Unmanned weapons (traps).
4.10.3. Non-compliant double ended weapons (e.g. nunchaku, double-ended daggers, pommel spikes).
4.10.4. Punching weapons (punching daggers, tonfas).
4.10.5. Any weapon that, when used as intended, violates the rules stipulated in the Book of War or grants an excessive advantage.

4.11. Shields:

4.11.1. Shields must be padded on the edges and face so as not to cause injury when struck with a forceful blow of an arm/hand.
4.11.2. The maximum width of a shield is three (3) feet (91.44 cm). Concave/curved shields will be measured along the curve of the face.
4.11.3. The maximum height of a shield is eighteen (18) inches (45.72 cm) less than the height of the wielder.
4.11.4. The minimum dimension on the face of a shield is twelve (12) inches (30.48 cm).
4.11.5. Shield spikes are allowed for decoration but may not form any rigid protrusions.
4.11.6. Shields must be reasonably rigid which is defined as the edges not bending towards each other excessively when attempting to bend the shield in half.

4.12. Armor Checking:

4.12.1. Definitions:

4.12.1.1. Composite: Armor of metal, leather, or both that is attached to another material backing and/or covering.
4.12.1.2. Cops: Rigid knee and elbow armor.
4.12.1.3. Gauntlet: Armor for the hands.
4.12.1.4. Gorget: Armor specifically for the neck.
4.12.1.5. Helmet: Armor for the head and neck.
4.12.1.6. Leather: Armor constructed of tanned animal hide. Synthetic leather, pleather, or vinyl, and other man-made materials cannot be used in place of actual animal hides.
4.12.1.7. Metal: Armor constructed of metal. Includes chain and plate.

4.12.1.7.1. Rigid Metal: Armor constructed of discrete or continuous metal plates.
4.12.1.7.2. Chain or Maille (mail): Metal Armor constructed of interlocking metal rings.

4.12.1.8. Penny round: Armor checking standard where the edge of rigid metal in armor is compared to that of a penny:

4.12.1.8.1. The edge of rigid metal armor shall have the smoothness of the edge of a penny.
4.12.1.8.2. The edge of rigid metal armor shall have less cutting ability than the edge of a penny.
4.12.1.8.3. The radius of any rigid metal corner must be greater than the radius of a penny.

4.12.1.9. Sabaton: Armor for the foot.

4.12.2. Armor must be inspected for safety by marshals.
4.12.3. Armor must not catch appendages such as fingers. This includes articulated plates and large diameter chain.
4.12.4. Armor may not have rigid protrusions that rise more than one-half (½) inch (1.27 cm) from the surface.
4.12.5. Leather Armor:

4.12.5.1. The minimum thickness for leather armor is 10oz (5/32'', 4.0 mm).).
4.12.5.2. The minimum thickness requirement can be achieved by layering up to two (2) pieces of thinner leather.

4.12.6. Metal Armor:

4.12.6.1. Metal Armor must be made from period metals and alloys such as iron, bronze, brass, or copper. Modern steel alloys are also allowed.
4.12.6.2. Metal Armor must conform to both of the following:

4.12.6.2.1. Must not be easily deformable by hand or by weapon strikes.
4.12.6.2.2. Using a material with a thickness of at least twenty (20) gauge.

4.12.6.3. Rigid metal must conform to the penny round standard.

4.12.7. Composite Armor:

4.12.7.1. Composite armor is defined as armor made up of leather, metal, or both that is attached to and/or covered by another material such as thin leather or cloth.
4.12.7.2. Studded, scaled, or brigandine armor can only be counted as Armor if two-thirds (2/3) of the armor piece is constructed from armor-grade metal or leather. The studs/rings/plates must be no more than one-half (½) inch (1.27 cm) apart in all directions. Rings and washers also cannot have openings larger than one-half (½) inch (1.27 cm).
4.12.7.3. Composite Armor must be readily identifiable as armor by appearance.

4.12.8. Prohibited Armor:

4.12.8.1. Rigid metal knee or elbow armor (cops).
4.12.8.2. Rigid metal hand armor.

5. GARB:
5.1. Garb is defined as the clothing to be worn on the Belegarth battlefield.
5.2. Minimum garb is the basic requirements for all participants. Minimum garb is defined as:

5.2.1. A tunic or tabard covering the torso.

5.2.1.1. Neutral colored t-shirts, with no visible printing, may be worn underneath a tunic or tabard.
5.2.1.2. Wearing nothing on the torso is acceptable, so long as it does not violate any applicable laws while at a BMCS event, including laws of local municipalities, states, provinces, territories, or federal governments.
5.2.1.3. Wearing neutral colored undergarments, base layers, or gender affirming wear (including but not limited to: sports bras, binders, or compression shirts) with no visible logos or modern prints is acceptable so long as it does not violate any applicable laws while at a BMCS event, including laws of local municipalities, states, provinces, territories, or federal governments.

5.2.2. Baggy pants or trousers covering the legs.
5.2.3. Skirts, Kilts, and Dresses are acceptable substitutes.
5.2.4. Footwear should be muted colors of a dark or natural color. Boots are preferred. Athletic shoes should be of a dark or natural color. Bare feet or Sandals are acceptable.

5.3. Modern clothing should be disguised, modified, or otherwise blend in with and be unobtrusive to the medieval or generally accepted Belegarth aesthetic.
5.4. Any piece of modern equipment or clothing required out of medical necessity overrules the minimum garb requirements. It should be covered when possible by clothing that does meet the minimum garb requirements or the item should be superficially altered accordingly.
5.5. Forbidden items:

5.5.1. T-shirts that are brightly colored, with visible logos, with visible collars, and or visible pockets.
5.5.2. Camouflage, cargo pants, or modern shorts.
5.5.3. Modern jeans of any color.
5.5.4. Modern hats.
5.5.5. Any fabrics with modern prints.
5.5.6. Any realistic weapons.
5.5.7. Cleats and spikes.

\end{document}